\begin{helper}
Let \(R\) and \(U\) be finite sets of tagged red and blue edge coordinates.
Assume that every coordinate in \(R\) and every coordinate in \(U\) belongs
to \(\noderef{TwoBiteTerminalCoordinateUniverse}(m)\), so each is a tagged
increasing non-loop base-edge coordinate, and assume \(U\cap R=\emptyset\).
Then every terminal truth table \(A\subseteq U\) extends any base sample
\(\omega_0\) to a genuine two-bite sample \(\omega_A\) with the same
embedding as \(\omega_0\), the same recorded coordinate values as
\(\omega_0\) on \(R\), and terminal coordinate values
\[
\noderef{TwoBiteEdgeCoordinateValue}(\omega_A,e)
  \quad\Longleftrightarrow\quad e\in A
\]
for every \(e\in U\).
\end{helper}
\begin{proof}
Unpack \(\noderef{TwoBiteSample}\): a sample is a red simple graph, a blue
simple graph, and an embedding.  Keep the embedding of \(\omega_0\) fixed.
It remains to define new red and blue simple graphs.

For each color separately, start with the corresponding graph of
\(\omega_0\).  For an unordered non-loop base edge whose tagged increasing
coordinate lies in \(U\), set its adjacency value to be the truth value
prescribed by \(A\).  For every other unordered base edge, keep the adjacency
value from \(\omega_0\).  This defines a symmetric loopless adjacency
relation: loops are always false, and for distinct endpoints the construction
first replaces the ordered pair by its unique increasing representative
before reading whether the tagged coordinate lies in \(U\) and \(A\).  Hence
the value assigned to \((x,y)\) is the same as the value assigned to
\((y,x)\).

The hypothesis \(U\subseteq\noderef{TwoBiteTerminalCoordinateUniverse}(m)\)
is exactly the orientation and loop hygiene needed in the previous paragraph:
no terminal coordinate asks us to set a loop, and no unordered base edge
appears twice with opposite orientations.  The same canonical-coordinate
hypothesis for \(R\), together with \(U\cap R=\emptyset\), says that no
recorded coordinate has the same tagged canonical representative as a
terminal coordinate being changed.  Therefore every recorded coordinate keeps
the value it had in \(\omega_0\).

For a terminal coordinate \(e\in U\), the new graph was defined precisely by
the membership test \(e\in A\), so
\(\noderef{TwoBiteEdgeCoordinateValue}(\omega_A,e)\) holds exactly when
\(e\in A\).  For a recorded coordinate \(c\in R\), the previous paragraph
shows that the graph adjacency read by
\(\noderef{TwoBiteEdgeCoordinateValue}\) is unchanged from \(\omega_0\).
Together with the unchanged embedding, these are the three required
properties of the extension sample.
\end{proof}
